Therefore, \eqref{2:connectW1a} gives
\begin{gather*} \lim_{z\to0^+} z^{-\mu-1} W_1(z,u)=\alpha_+(u)\end{gather*}
and, from \eqref{1:W1}, \eqref{1:W1a}, \eqref{1:As}
\begin{gather*} \lim_{z\to0^+} z^{-\mu-1}W_1(u,z)=\frac{2^{-\mu}u^\mu}{\Gamma(\mu+1)}\left(1+O\left(\frac1{u^{2N}}\right)\right)\end{gather*}
which implies \eqref{2:alpha} with $\alpha(u)=\alpha_+(u)$.
\end{proof}

Let us def\/ine
\begin{gather*}%\label{2:W3}
W_3(u,\mu,z)= \frac{2^{-\mu}u^\mu}{\Gamma(\mu+1)}W_+(u,\mu,z) .
\end{gather*}
Then Lemma~\ref{2:l1} gives
\begin{gather*}%\label{2:connectW3}
 W_3(u,z)=\tilde\alpha(u) W_1(u,z),\qquad \text{where} \quad \tilde\alpha(u)=1+O\left(\frac1{u^{2N}}\right).
\end{gather*}
Therefore, $W_3$ admits the asymptotic expansion \eqref{1:W1}, \eqref{1:W1a}, so we can replace $W_1$ by $W_3$.
Note that in contrast to $W_1$, $W_3$ is a uniquely def\/ined function which is identif\/ied as a (Floquet) solution of~\eqref{1:eq1} and not by its asymptotic behavior as $u\to\infty$.

Unfortunately, it seems impossible to replace $W_2$ by an easily identif\/iable solution of~\eqref{1:eq1}. However, we will now prove several useful properties of~$W_2$.

\begin{Lemma}\label{2:l2}
Suppose that $\Re\mu\ge 0$. There is a function $\beta(u)$ such that
\begin{gather}\label{2:eq3}
 W_2\big(u,ze^{\pi i}\big)-e^{\pi i(1-\mu)}W_2(u,z)= \beta(u) W_3(u,z) ,
\end{gather}
and, for every $N=1,2,3,\dots$,
\begin{gather}\label{2:beta}
 \beta(u)=\pi i\left(1+O\left(\frac{1}{u^{2N}}\right)\right)\qquad\text{as}\quad 0<u\to\infty.
 \end{gather}
\end{Lemma}
\begin{proof}
We set $\lambda_\pm=e^{\pi i(1\pm\mu)}$. Equation \eqref{1:eq1} has a fundamental system of solutions~$W_+$,~$W_-$ such that
\begin{gather*}
W_+\big(ze^{\pi i}\big)=\lambda_+ W_+(z),\qquad
W_-\big(ze^{\pi i}\big)=\lambda_- W_-(z)+\rho W_+(z) .
\end{gather*}
Let $w(z)=c_+W_+(z)+c_-W_-(z)$ be any solution of \eqref{1:eq1}.
Then
\begin{gather*}
w\big(ze^{\pi i}\big)-\lambda_- w(z)=((\lambda_+-\lambda_-)c_++\rho c_-) W_+(z).
\end{gather*}
If we apply this result to $w=W_2$ we see that there is a function~$\beta(u)$ such that~\eqref{2:eq3} holds.

Let $z>0$ and set $z_1=ze^{\pi i}$.
We use~\eqref{1:W2} for $z_1$ in place of~$z$, and \cite[(10.34.2)]{NIST}
\begin{gather}\label{3:ancontK}
K_\nu\big(ze^{\pi im}\big)= e^{-\pi i \nu m}K_\nu(z)-\pi i \frac{\sin(\pi \nu m)}{\sin(\pi \nu)} I_\nu(z)
\end{gather}
with $m=1$. Then
\begin{gather*}
 W_2(u,z_1) = z\left(\lambda_-K_\mu(uz)+\pi i I_\mu(uz)\right)\left(\sum_{s=0}^{N-1}\frac{A_s(z)}{u^{2s}}+g_2(u,z_1)\right)\\
\hphantom{W_2(u,z_1) =}{} +\frac{z}{u}\left(-\lambda_-K_{\mu+1}(uz)+\pi i I_{\mu+1}(uz)\right) \left(\sum_{s=0}^{N-1}\frac{B_s(z)}{u^{2s}}+zh_2(u,z_1)\right).
\end{gather*}
Using \eqref{1:W2} a second time, we f\/ind that
\begin{gather*}
 W_2(u,z_1)-\lambda_-W_2(u,z)=\pi i zI_\mu(uz)\left(\sum_{s=0}^{N-1}\frac{A_s(z)}{u^{2s}}+g_2(u,z_1)\right)\\
\qquad{}
+\pi i\frac{z}{u}I_{\mu+1}(uz) \left(\sum_{s=0}^{N-1}\frac{B_s(z)}{u^{2s}}+zh_2(u,z_1)\right)\\
\qquad{}
+\lambda_- zK_\mu(uz) (g_2(u,z_1)-g_2(u,z))-\lambda_-\frac{z^2}{u}K_{\mu+1}(uz)(h_2(u,z_1)-h_2(u,z)).
\end{gather*}
We now expand the right-hand side of \eqref{2:eq3} using \eqref{1:W1}, and compare the expansions.
Setting $z=R$ and dividing by $R I_\mu(uR)$, we obtain
\begin{gather*}
 (\beta(u)-\pi i) \left(1+O\left(\frac1u\right)\right)=O\left(\frac{1}{u^{2N}}\right) \quad \text{as $0<u\to\infty$,}
\end{gather*}
where we used \cite[(10.40.1)]{NIST}
\begin{gather}\label{2:asyI}
 I_\nu(x)=\frac{e^x}{\sqrt{2\pi x}}\left(1+O\left(\frac1x\right)\right)\qquad\text{as} \quad 0<x\to\infty,
\end{gather}
and \cite[(10.40.2)]{NIST}
\begin{gather}\label{2:asyK}
 K_\nu(x)=\sqrt{\frac\pi{2x}}e^{-x}\left(1+O\left(\frac1x\right)\right)\qquad\text{as} \quad 0<x\to\infty.
\end{gather}
This proves \eqref{2:beta}.
\end{proof}
